\ticket{11}{Методы поиска решений в пространстве состояний. Метод ветвей и границ (Branch and Bound), лучевой поиск (Beam Search), А*-алгоритм и его частные случаи}

\begin{examplan}
    \item Описать задачу поиска в пространстве состояний.
    \item Объяснить метод ветвей и границ.
    \item Объяснить лучевой поиск.
    \item Описать A*-алгоритм, условия оптимальности и частные случаи.
\end{examplan}

\subsection*{Короткая версия для письменного ответа}

\textbf{Пространство состояний} --- граф, где вершины соответствуют состояниям задачи, а дуги --- действиям, переводящим одно состояние в другое. Цель поиска --- найти путь от начального состояния к целевому, часто с минимальной стоимостью.

Основные компоненты: начальное состояние, операторы, целевой тест, стоимость пути и стратегия выбора следующего состояния для раскрытия.

\subsection*{Метод ветвей и границ}

\textbf{Branch and Bound} --- метод оптимизации, который строит дерево вариантов и отсекает ветви, заведомо не способные улучшить уже найденное решение.

Основные шаги:
\begin{enumerate}
    \item \textbf{Ветвление}: задача разбивается на подзадачи.
    \item \textbf{Оценка границ}: для каждой подзадачи вычисляется нижняя граница для минимизации или верхняя для максимизации.
    \item \textbf{Отсечение}: если граница хуже текущего рекорда, ветвь не рассматривается.
    \item \textbf{Обновление рекорда}: найденное допустимое решение сравнивается с лучшим известным.
\end{enumerate}

Качество метода зависит от функции оценки границ и стратегии выбора узла: в глубину, в ширину или по наилучшей оценке. Применения: задача коммивояжера, целочисленное программирование, задача назначений.

\subsection*{Лучевой поиск}

\textbf{Beam Search} --- эвристический поиск, который на каждом уровне оставляет только $B$ лучших кандидатов по оценочной функции. Число $B$ называют шириной луча.

Алгоритм:
\begin{enumerate}
    \item сгенерировать преемников текущих узлов;
    \item оценить их эвристикой;
    \item оставить только $B$ лучших;
    \item повторять до цели или исчерпания вариантов.
\end{enumerate}

Лучевой поиск экономит память и время, но не является полным и не гарантирует оптимальность: хорошая ветвь может быть отброшена слишком рано. При $B=1$ он близок к жадному поиску.

\subsection*{A*-алгоритм}

\textbf{A*} --- информированный алгоритм поиска кратчайшего пути, выбирающий узел с минимальной оценкой:
\[
f(n)=g(n)+h(n),
\]
где $g(n)$ --- известная стоимость пути от старта до $n$, а $h(n)$ --- эвристическая оценка стоимости от $n$ до цели.

Шаги A*:
\begin{enumerate}
    \item поместить стартовый узел в открытый список;
    \item выбрать узел с минимальным $f(n)$;
    \item если это цель, восстановить путь;
    \item иначе раскрыть его преемников;
    \item для каждого преемника обновить $g$, $h$, $f$ и родителя, если найден лучший путь;
    \item повторять, пока цель не найдена или открытый список не пуст.
\end{enumerate}

\subsection*{Свойства A*}

\begin{itemize}
    \item \textbf{Допустимая эвристика}: $0\leq h(n)\leq h^*(n)$, то есть эвристика не переоценивает реальную стоимость до цели.
    \item \textbf{Согласованная эвристика}: $h(n)\leq c(n,n')+h(n')$ для каждого перехода $n\to n'$.
    \item При допустимой эвристике A* оптимален для деревьев поиска; при согласованной --- удобно работает на графах без повторного раскрытия закрытых узлов.
    \item Полнота выполняется при положительных стоимостях шагов и конечном ветвлении.
\end{itemize}

\subsection*{Частные случаи}

\begin{description}
    \item[Алгоритм Дейкстры] получается при $h(n)=0$: выбор идет только по уже накопленной стоимости $g(n)$.
    \item[Поиск в ширину] получается при $h(n)=0$ и единичных стоимостях всех шагов.
    \item[Жадный best-first search] получается, если игнорировать $g(n)$ и выбирать по $h(n)$; он часто быстр, но не гарантирует оптимальность.
\end{description}

\begin{important}
    \item Задача поиска задается начальными состоянием, операторами, целью и стоимостью.
    \item Branch and Bound ищет оптимум и отсекает ветви по оценкам границ.
    \item Beam Search ограничивает число рассматриваемых кандидатов и жертвует полнотой.
    \item A* балансирует уже пройденную стоимость $g$ и прогноз $h$.
    \item Допустимость эвристики нужна для оптимальности A*.
\end{important}
